\documentclass[../../main.tex]{subfiles}
\begin{document}
\section{多重积分}\label{section:appC-multiple-integrals}

本节中我们简要回顾 $\mathbb R^n$ 中多重积分的一些基本事实。对我们的目的，Riemann 积分已绰绰有余。熟悉 Lebesgue 积分理论的读者可自由地在 Lebesgue 意义下解释我们的所有积分，因为对所考虑的函数类型，这两种积分相等。关于这里所述积分论方面的更多细节，可查阅 [Apo74, Rud76, Str00]。

$\mathbb R^n$ 中的\emph{闭矩形}是形如 $[a^1,b^1]\times\cdots\times[a^n,b^n]$（其中实数 $a^i<b^i$）的乘积集合。类似地，\emph{开矩形}是形如 $(a^1,b^1)\times\cdots\times(a^n,b^n)$ 的集合。若 $A$ 是任一类型的矩形，$A$ 的\emph{体积}记作 $\operatorname{Vol}(A)$，定义为其各分量区间长度之积：
\begin{align}\label{align:appC-rectangle-volume}
\operatorname{Vol}(A)=(b^1-a^1)\cdots(b^n-a^n).
\end{align}
如果一个矩形的所有边长 $(b^i-a^i)$ 都相等，就称为\emph{立方体}。

给定闭区间 $[a,b]\subset\mathbb R$，$[a,b]$ 的一个\emph{划分}是满足 $a=a_0<a_1<\cdots<a_k=b$ 的实数列 $P=(a_0,\dots,a_k)$。每个区间 $[a_{i-1},a_i]$（$i=1,\dots,k$）称为 $P$ 的一个\emph{子区间}。类似地，如果 $A=[a^1,b^1]\times\cdots\times[a^n,b^n]$ 是闭矩形，$A$ 的一个\emph{划分}是一个 $n$ 元组 $P=(P_1,\dots,P_n)$，其中每个 $P_i$ 是 $[a^i,b^i]$ 的划分。每个形如 $I_1\times\cdots\times I_n$（其中 $I_j$ 是 $P_j$ 的子区间）的矩形称为 $P$ 的一个\emph{子矩形}。显然，$A$ 是任一划分中所有子矩形之并，而不同子矩形只在边界上相交。

设 $A\subset\mathbb R^n$ 是闭矩形，$f:A\to\mathbb R$ 是有界函数。对 $A$ 的每个划分 $P$，定义 $f$ 关于 $P$ 的\emph{下和}为
\begin{align*}
L(f;P)=\sum_j\left(\inf_{R_j}f\right)\operatorname{Vol}(R_j),
\end{align*}
其中求和遍历 $P$ 的所有子矩形 $R_j$。类似地，\emph{上和}为
\begin{align*}
U(f;P)=\sum_j\left(\sup_{R_j}f\right)\operatorname{Vol}(R_j).
\end{align*}
关于 $P$ 的下和显然小于或等于关于同一划分的上和。事实上还有更强的结论。

\begin{lemma}\label{lemma:appC-upper-lower}
设 $A\subset\mathbb R^n$ 是闭矩形，$f:A\to\mathbb R$ 是有界函数。对 $A$ 的任意一对划分 $P$ 与 $P'$，
\begin{align*}
L(f;P)\leqslant U(f;P').
\end{align*}
\end{lemma}

\begin{proof}
写 $P=(P_1,\dots,P_n)$、$P'=(P'_1,\dots,P'_n)$，令 $Q$ 为划分 $Q=(P_1\cup P'_1,\dots,P_n\cup P'_n)$。$P$ 或 $P'$ 的每个子矩形都是 $Q$ 的有限个子矩形之并。简单计算表明
\begin{align*}
L(f;P)\leqslant L(f;Q)\leqslant U(f;Q)\leqslant U(f;P'),
\end{align*}
由此得结果。
\end{proof}

$f$ 在 $A$ 上的\emph{下积分}为
\begin{align*}
\underline{\int_A}f\,\mathrm{d}V=\sup\{L(f;P):P\text{ 是 }A\text{ 的划分}\},
\end{align*}
而\emph{上积分}为
\begin{align*}
\overline{\int_A}f\,\mathrm{d}V=\inf\{U(f;P):P\text{ 是 }A\text{ 的划分}\}.
\end{align*}
显然，由于 $f$ 有界，这两个数都存在；引理 C.17 蕴含下积分小于或等于上积分。

如果有界函数 $f:A\to\mathbb R$ 的上积分与下积分相等，就称 $f$ 在 $A$ 上（Riemann）\emph{可积}，它们的公共值记作 $\int_A f\,\mathrm{d}V$，称为 $f$ 在 $A$ 上的\emph{积分}。此记号中的“$\mathrm{d}V$”，正如单积分记号中的“$\mathrm{d}x$”，本身没有意义；它只是积分号的一个“闭括号”。其他常见记号有 $\int_A f$、$\int_A f\,\mathrm{d}x^1\cdots\mathrm{d}x^n$ 或 $\int_A f(x^1,\dots,x^n)\,\mathrm{d}x^1\cdots\mathrm{d}x^n$。在 $\mathbb R^2$ 中，符号 $\mathrm{d}V$ 常换成 $\mathrm{d}A$。

有界函数 Riemann 可积有一个简单判别法。它基于下列概念。如果对每个 $\delta>0$，存在 $X$ 的一个开矩形可数覆盖 $\{C_i\}$，使 $\sum_i\operatorname{Vol}(C_i)<\delta$，就称子集 $X\subset\mathbb R^n$ 有\emph{零测度}。（熟悉 Lebesgue 测度理论的读者会注意到，这等价于 $X$ 的 Lebesgue 测度等于零。）

\begin{proposition}[零测集的性质]\label{proposition:appC-null-set-properties}
\begin{enumerate}[(a)]
\item 如果 $X\subset\mathbb R^n$ 有零测度且 $x_0\in\mathbb R^n$，则平移子集 $x_0+X=\{x_0+a:a\in X\}$ 也有零测度。
\item $\mathbb R^n$ 中零测集的每个子集都有零测度。
\item $\mathbb R^n$ 中零测集的可数并仍有零测度。
\item 如果 $k<n$，则 $\mathbb R^k$ 的每个子集（看作 $\mathbb R^n$ 中满足 $x^{k+1}=\cdots=x^n=0$ 的点 $x$ 所成集合）在 $\mathbb R^n$ 中有零测度。
\end{enumerate}
\end{proposition}

\begin{exercise}\label{exercise:appC-C-19}
证明命题 C.18。
\end{exercise}

本命题的 (d) 说明：有零测度是一个集合相对于它所在的某个 Euclid 空间的性质，而非集合本身固有的性质。例如，$x$ 轴上的一个开区间作为 $\mathbb R^2$ 的子集有零测度，但作为 $\mathbb R^1$ 的子集则不然。因此，当我们想强调某子集作为 $\mathbb R^n$ 的子集有零测度时，有时说它有 $n$ 维零测度。

下列命题给出函数可积的一个充分条件。它特别表明每个有界连续函数都可积。

\begin{proposition}[Lebesgue 可积判别法]\label{proposition:appC-lebesgue-integrability}
设 $A\subset\mathbb R^n$ 是闭矩形，$f:A\to\mathbb R$ 是有界函数。如果集合
\begin{align*}
S=\{x\in A:f\text{ 在 }x\text{ 处不连续}\}
\end{align*}
有零测度，则 $f$ 可积。
\end{proposition}

\begin{proof}
假设集合 $S$ 有零测度，给定 $\varepsilon>0$。由零测集的定义，$S$ 可被开矩形可数族 $\{C_i\}$ 覆盖，其体积之和小于 $\varepsilon$。对每个 $q\in A\setminus S$，由于 $f$ 在 $q$ 处连续，存在以 $q$ 为中心的开矩形 $D_q$，使对所有 $x\in D_q\cap A$ 有 $|f(x)-f(q)|<\varepsilon$。把 $D_q$ 稍稍缩小，可安排同样的不等式对所有 $x\in\overline{D_q}\cap A$ 成立。这蕴含 $\sup_{\overline{D_q}}f-\inf_{\overline{D_q}}f\leqslant2\varepsilon$。

所有形如 $C_i$ 或 $D_q$ 的矩形所成族是 $A$ 的一个开覆盖。由紧性，其中有限多个覆盖 $A$。把这些矩形重新标记为 $\{C_1,\dots,C_k,D_1,\dots,D_l\}$。把每个 $C_i$ 或 $D_j$ 换成它与 $\operatorname{Int}A$ 的交，可假定每个 $\overline{C_i}$ 和每个 $\overline{D_j}$ 都含于 $A$。

由于矩形 $\{C_i,D_j\}$ 只有有限多个，存在 $A$ 的划分 $P$，使每个 $\overline{C_i}$ 或 $\overline{D_j}$ 都等于 $P$ 的若干子矩形之并。（只需用矩形 $C_i$ 和 $D_j$ 各分量区间的所有端点之并来定义划分。）我们可把 $P$ 的子矩形分成两个不相交集合 $\mathcal C$ 与 $\mathcal D$，使 $\mathcal C$ 中每个子矩形都含于某个 $\overline{C_i}$，$\mathcal D$ 中每个子矩形都含于某个 $\overline{D_j}$。则
\begin{align*}
U(f;P)-L(f;P)
&=\sum_i\left(\sup_{R_i}f\right)\operatorname{Vol}(R_i)-\sum_i\left(\inf_{R_i}f\right)\operatorname{Vol}(R_i)\\
&=\sum_{R_i\in\mathcal C}\left(\sup_{R_i}f-\inf_{R_i}f\right)\operatorname{Vol}(R_i)
+\sum_{R_i\in\mathcal D}\left(\sup_{R_i}f-\inf_{R_i}f\right)\operatorname{Vol}(R_i)\\
&\leqslant\left(\sup_A f-\inf_A f\right)\sum_{R_i\in\mathcal C}\operatorname{Vol}(R_i)
+2\varepsilon\sum_{R_i\in\mathcal D}\operatorname{Vol}(R_i)\\
&\leqslant\left(\sup_A f-\inf_A f\right)\varepsilon+2\varepsilon\operatorname{Vol}(A).
\end{align*}
由此
\begin{align*}
\overline{\int_A}f\,\mathrm{d}V-\underline{\int_A}f\,\mathrm{d}V
\leqslant\left(\sup_A f-\inf_A f\right)\varepsilon+2\varepsilon\operatorname{Vol}(A),
\end{align*}
取 $\varepsilon$ 足够小即可使之任意小。这蕴含 $f$ 的上积分与下积分相等，故 $f$ 可积。
\end{proof}

事实上，Lebesgue 判别法对 Riemann 可积既是必要的也是充分的，但我们不需要它。

现设 $D\subset\mathbb R^n$ 是任意有界集，$f:D\to\mathbb R$ 是有界函数。用
\begin{align}\label{align:appC-extended-function}
f_D(x)=\begin{cases}f(x),&x\in D;\\0,&x\in\mathbb R^n\setminus D\end{cases}
\end{align}
定义 $f_D:\mathbb R^n\to\mathbb R$。如果对某个含 $D$ 的闭矩形 $A$，积分
\begin{align}\label{align:appC-integral-extended}
\int_A f_D\,\mathrm{d}V
\end{align}
存在，就称 $f$ 在 $D$ 上\emph{可积}。积分\eqref{align:appC-integral-extended} 记作 $\int_D f\,\mathrm{d}V$，称为 $f$ 在 $D$ 上的积分。容易验证：$f$ 的可积性与积分值都与所取矩形无关。

实际上，我们只关心有界连续函数的积分。然而，由于我们有时需要在矩形以外的区域上积分它们，也有必要考虑不连续函数的积分，例如由\eqref{align:appC-extended-function} 定义的函数 $f_D$。证明命题 C.20 的主要原因是：它使我们能简单描述那种所有有界连续函数都可积的区域。

如果 $D$ 有界且 $\partial D$ 有 $n$ 维零测度，就称子集 $D\subset\mathbb R^n$ 为\emph{积分区域}。由命题 C.18 可得：每个开矩形或闭矩形都是积分区域，有限个积分区域之并仍是积分区域。

\begin{proposition}\label{proposition:appC-integral-domain}
如果 $D\subset\mathbb R^n$ 是积分区域，则 $D$ 上每个有界连续实值函数都在 $D$ 上可积。
\end{proposition}

\begin{proof}
设 $f:D\to\mathbb R$ 有界且连续，$f_D:\mathbb R^n\to\mathbb R$ 是由\eqref{align:appC-extended-function} 定义的函数，$A$ 是含 $D$ 的闭矩形。为证定理，只需证明 $A$ 中 $f_D$ 不连续的点集有零测度。

如果 $x\in\operatorname{Int}D$，则 $f_D=f$ 在 $x$ 的某邻域上成立，故 $f_D$ 在 $x$ 处连续。类似地，如果 $x\in\mathbb R^n\setminus\overline D$，则 $f_D\equiv0$ 在 $x$ 的某邻域上成立，故 $f$ 又在 $x$ 处连续。因而 $f_D$ 不连续的点集含于 $\partial D$，故有零测度。
\end{proof}

当然，如果 $D$ 紧，则前一命题中 $f$ 有界的假设是多余的。

如果 $D$ 是积分区域，$D$ 的\emph{体积}定义为
\begin{align}\label{align:appC-domain-volume}
\operatorname{Vol}(D)=\int_D 1\,\mathrm{d}V.
\end{align}
右端积分常缩写为 $\int_D\mathrm{d}V$。

下列两个命题汇集关于连续函数体积与积分的一些基本事实。

\begin{proposition}[体积的性质]\label{proposition:appC-volume-properties}
设 $D\subset\mathbb R^n$ 是积分区域。
\begin{enumerate}[(a)]
\item 如果 $D$ 是开矩形或闭矩形，则 $\operatorname{Vol}(D)$ 的两个定义\eqref{align:appC-rectangle-volume} 与\eqref{align:appC-domain-volume} 一致。
\item $\operatorname{Vol}(D)\geqslant0$，等号当且仅当 $D$ 有零测度。
\item 如果 $D_1,\dots,D_k$ 是积分区域且其并为 $D$，则
\begin{align*}
\operatorname{Vol}(D)\leqslant\operatorname{Vol}(D_1)+\cdots+\operatorname{Vol}(D_k),
\end{align*}
等号当且仅当对每个 $i\neq j$，$D_i\cap D_j$ 有零测度。
\item 如果 $D_1$ 是含于 $D$ 的积分区域，则 $\operatorname{Vol}(D_1)\leqslant\operatorname{Vol}(D)$，等号当且仅当 $D\setminus D_1$ 有零测度。
\end{enumerate}
\end{proposition}

\begin{proposition}[积分的性质]\label{proposition:appC-integral-properties}
设 $D\subset\mathbb R^n$ 是积分区域，$f,g:D\to\mathbb R$ 连续且有界。
\begin{enumerate}[(a)]
\item 对任意 $a,b\in\mathbb R$，
\begin{align*}
\int_D(af+bg)\,\mathrm{d}V=a\int_D f\,\mathrm{d}V+b\int_D g\,\mathrm{d}V.
\end{align*}
\item 如果 $D$ 有零测度，则 $\int_D f\,\mathrm{d}V=0$。
\item 如果 $D_1,\dots,D_k$ 是积分区域，其并为 $D$，且其两两交有零测度，则
\begin{align*}
\int_D f\,\mathrm{d}V=\int_{D_1}f\,\mathrm{d}V+\cdots+\int_{D_k}f\,\mathrm{d}V.
\end{align*}
\item 如果 $f\geqslant0$ 在 $D$ 上，则 $\int_D f\,\mathrm{d}V\geqslant0$，等号当且仅当 $f\equiv0$ 在 $\operatorname{Int}D$ 上。
\item $(\inf_D f)\operatorname{Vol}(D)\leqslant\int_D f\,\mathrm{d}V\leqslant(\sup_D f)\operatorname{Vol}(D)$。
\item $\left|\int_D f\,\mathrm{d}V\right|\leqslant\int_D|f|\,\mathrm{d}V$。
\end{enumerate}
\end{proposition}

\begin{exercise}\label{exercise:appC-C-24}
证明命题 C.22 与 C.23。
\end{exercise}

\begin{corollary}\label{corollary:appC-no-open-null}
$\mathbb R^n$ 中的零测集不含非空开子集。
\end{corollary}

\begin{proof}
为反证，假设 $D\subset\mathbb R^n$ 有零测度且含非空开子集 $U$。则 $U$ 含非空开矩形，它有正体积，故由命题 C.22 没有零测度。但由命题 C.18，$D$ 的每个子集都有零测度，矛盾。
\end{proof}

我们还需要多重积分的另外两个基本性质。证明过于繁复，不能收入本综述，但如果你感兴趣可在 [Apo74, Rud76, Str00] 中查到。这些定理每种都可有多种陈述方式，有的比其他更强。这里给出的版本对我们的应用已足够。

\begin{theorem}[变量替换]\label{theorem:appC-change-of-variables}
假设 $D,E$ 是 $\mathbb R^n$ 中开的积分区域，$G:\overline D\to\overline E$ 是光滑映射，且限制为从 $D$ 到 $E$ 的微分同胚。对每个连续函数 $f:\overline E\to\mathbb R$，
\begin{align*}
\int_E f\,\mathrm{d}V=\int_D(f\circ G)|\det \mathrm{D}G|\,\mathrm{d}V.
\end{align*}
\end{theorem}

\begin{theorem}[Fubini 定理]\label{theorem:appC-fubini}
设 $A=[a^1,b^1]\times\cdots\times[a^n,b^n]$ 是 $\mathbb R^n$ 中的闭矩形，$f:A\to\mathbb R$ 连续。则
\begin{align*}
\int_A f\,\mathrm{d}V
=\int_{a^n}^{b^n}\left(\cdots\int_{a^1}^{b^1}f(x^1,\dots,x^n)\,\mathrm{d}x^1\cdots\right)\mathrm{d}x^n,
\end{align*}
且右端累次积分中的变量以任意方式重排也成立。
\end{theorem}

\subsection*{向量值函数的积分}

如果 $D\subset\mathbb R^n$ 是积分区域，$F:D\to\mathbb R^k$ 是有界连续向量值函数，我们定义 $F$ 在 $D$ 上的积分为把 $F$ 逐分量积分所得的 $\mathbb R^k$ 中向量：
\begin{align*}
\int_D F\,\mathrm{d}V=\left(\int_D F^1\,\mathrm{d}V,\dots,\int_D F^k\,\mathrm{d}V\right).
\end{align*}
命题 C.23 的 (a)--(c) 对向量值积分显然也成立，只需逐分量应用。(f) 也成立，但证明需稍多功夫。

\begin{proposition}\label{proposition:appC-vector-integral-bound}
假设 $D\subset\mathbb R^n$ 是积分区域，$F:D\to\mathbb R^k$ 是有界连续向量值函数。则
\begin{align}\label{align:appC-vector-integral-bound}
\left|\int_D F\,\mathrm{d}V\right|\leqslant\int_D|F|\,\mathrm{d}V.
\end{align}
\end{proposition}

\begin{proof}
令 $G$ 表示向量 $\int_D F\,\mathrm{d}V\in\mathbb R^k$。如果 $G=0$，则\eqref{align:appC-vector-integral-bound} 显然成立，故不妨设 $G\neq0$。计算得
\begin{align*}
|G|^2=\sum_{i=1}^k(G^i)^2=\sum_{i=1}^k G^i\int_D F^i\,\mathrm{d}V
=\sum_{i=1}^k\int_D G^iF^i\,\mathrm{d}V=\int_D(G\cdot F)\,\mathrm{d}V.
\end{align*}
把命题 C.23(f) 用于数量积分 $\int_D(G\cdot F)\,\mathrm{d}V$，得
\begin{align*}
|G|^2\leqslant\int_D|G\cdot F|\,\mathrm{d}V\leqslant\int_D|G||F|\,\mathrm{d}V
=|G|\int_D|F|\,\mathrm{d}V.
\end{align*}
上式两边除以 $|G|$ 即得\eqref{align:appC-vector-integral-bound}。
\end{proof}

作为\eqref{align:appC-vector-integral-bound} 的应用，我们证明 $C^1$ 函数局部行为用其全导数表达的一个重要估计。

\begin{proposition}[$C^1$ 函数的 Lipschitz 估计]\label{proposition:appC-lipschitz}
设 $U\subset\mathbb R^n$ 是开子集，假设 $F:U\to\mathbb R^m$ 是 $C^1$ 类的。则 $F$ 在每个紧凸子集 $K\subset U$ 上 Lipschitz 连续。Lipschitz 常数可取为 $\sup_{x\in K}|\mathrm{D}F(x)|$。
\end{proposition}

\begin{proof}
由于 $|\mathrm{D}F(x)|$ 是 $x$ 的连续函数，它在紧集 $K$ 上有界。（这里范数是\eqref{align:appB-frobenius} 中定义的矩阵 Frobenius 范数。）令 $M=\sup_{x\in K}|\mathrm{D}F(x)|$。对任意 $a,b\in K$，由于 $K$ 凸，对所有 $t\in I$ 有 $a+t(b-a)\in K$。由对 $F$ 每个分量应用的微积分基本定理连同链式法则，
\begin{align*}
F(b)-F(a)=\int_0^1\frac{\mathrm{d}}{\mathrm{d}t}F(a+t(b-a))\,\mathrm{d}t
=\int_0^1\mathrm{D}F(a+t(b-a))(b-a)\,\mathrm{d}t.
\end{align*}
因此由\eqref{align:appC-vector-integral-bound} 和习题 B.48，
\begin{align*}
|F(b)-F(a)|\leqslant\int_0^1|\mathrm{D}F(a+t(b-a))|\,|b-a|\,\mathrm{d}t
\leqslant\int_0^1 M|b-a|\,\mathrm{d}t=M|b-a|.
\end{align*}
\end{proof}

\begin{corollary}\label{corollary:appC-local-lipschitz}
如果 $U\subset\mathbb R^n$ 是开子集，$F:U\to\mathbb R^m$ 是 $C^1$ 类的，则 $F$ 局部 Lipschitz 连续。
\end{corollary}

\begin{proof}
$U$ 的每一点都含于一个闭包含于 $U$ 的球中，命题 C.29 表明 $F$ 在这种球上的限制 Lipschitz 连续。
\end{proof}
\end{document}
